\documentclass[10pt,a4paper]{amsart}
\usepackage{amsmath,amssymb,amsthm,graphicx}
\usepackage{verbatim}
\newenvironment{reference}{\comment}{\endcomment}
%\newenvironment{reference}{}{}
\newenvironment{slogan}{\comment}{\endcomment}
\newenvironment{history}{\comment}{\endcomment}

% For commutative diagrams we use Xy-pic
\usepackage[all]{xy}

% We use 2cell for 2-commutative diagrams.
\xyoption{2cell}
% Theorem environments.
%
\theoremstyle{plain}
\newtheorem{theorem}[subsection]{Theorem}
\newtheorem{proposition}[subsection]{Proposition}
\newtheorem{lemma}[subsection]{Lemma}

\theoremstyle{definition}
\newtheorem{definition}[subsection]{Definition}
\newtheorem{example}[subsection]{Example}
\newtheorem{exercise}[subsection]{Exercise}
\newtheorem{situation}[subsection]{Situation}

\theoremstyle{remark}
\newtheorem{remark}[subsection]{Remark}
\newtheorem{remarks}[subsection]{Remarks}

\numberwithin{equation}{subsection}

% Macros
%
\def\lim{\mathop{\mathrm{lim}}\nolimits}
\def\colim{\mathop{\mathrm{colim}}\nolimits}
\def\Spec{\mathop{\mathrm{Spec}}}
\def\Hom{\mathop{\mathrm{Hom}}\nolimits}
\def\Ext{\mathop{\mathrm{Ext}}\nolimits}
\def\SheafHom{\mathop{\mathcal{H}\!\mathit{om}}\nolimits}
\def\SheafExt{\mathop{\mathcal{E}\!\mathit{xt}}\nolimits}
\def\Sch{\mathit{Sch}}
\def\Mor{\mathop{\mathrm{Mor}}\nolimits}
\def\Ob{\mathop{\mathrm{Ob}}\nolimits}
\def\Sh{\mathop{\mathit{Sh}}\nolimits}
\def\NL{\mathop{N\!L}\nolimits}
\def\CH{\mathop{\mathrm{CH}}\nolimits}
\def\proetale{{pro\text{-}\acute{e}tale}}
\def\etale{{\acute{e}tale}}
\def\QCoh{\mathit{QCoh}}
\def\Ker{\mathop{\mathrm{Ker}}}
\def\Im{\mathop{\mathrm{Im}}}
\def\Coker{\mathop{\mathrm{Coker}}}
\def\Coim{\mathop{\mathrm{Coim}}}

% Boxtimes
%
\DeclareMathSymbol{\boxtimes}{\mathbin}{AMSa}{"02}

%
% Macros for moduli stacks/spaces
%
\def\QCohstack{\mathcal{QC}\!\mathit{oh}}
\def\Cohstack{\mathcal{C}\!\mathit{oh}}
\def\Spacesstack{\mathcal{S}\!\mathit{paces}}
\def\Quotfunctor{\mathrm{Quot}}
\def\Hilbfunctor{\mathrm{Hilb}}
\def\Curvesstack{\mathcal{C}\!\mathit{urves}}
\def\Polarizedstack{\mathcal{P}\!\mathit{olarized}}
\def\Complexesstack{\mathcal{C}\!\mathit{omplexes}}
% \Pic is the operator that assigns to X its picard group, usage \Pic(X)
% \Picardstack_{X/B} denotes the Picard stack of X over B
% \Picardfunctor_{X/B} denotes the Picard functor of X over B
\def\Pic{\mathop{\mathrm{Pic}}\nolimits}
\def\Picardstack{\mathcal{P}\!\mathit{ic}}
\def\Picardfunctor{\mathrm{Pic}}
\def\Deformationcategory{\mathcal{D}\!\mathit{ef}}
\usepackage{cleveref}
\crefname{originallemma}{original lemma}{original lemmas}
\Crefname{originallemma}{Original lemma}{Original lemmas}
\makeatletter
\newlabel{divisors-lemma-blowing-up-projective}{{02NS}{source}{}{sourcepointer-1}{}}
\newlabel{divisors-lemma-blowing-up-projective@cref}{{[originallemma][0][]02NS}{[1][0][]source}{}{}{}}
\makeatother
\crefname{originalreference}{original reference}{original references}
\Crefname{originalreference}{Original reference}{Original references}
\makeatletter
\newlabel{lemma-cohomology-of-blowup}{{0AGS}{source}{}{sourcepointer-2}{}}
\newlabel{lemma-cohomology-of-blowup@cref}{{[originalreference][0][]0AGS}{[1][0][]source}{}{}{}}
\makeatother
\crefname{originallemma}{original lemma}{original lemmas}
\Crefname{originallemma}{Original lemma}{Original lemmas}
\makeatletter
\newlabel{morphisms-lemma-locally-projective-proper}{{01WC}{source}{}{sourcepointer-3}{}}
\newlabel{morphisms-lemma-locally-projective-proper@cref}{{[originallemma][0][]01WC}{[1][0][]source}{}{}{}}
\makeatother
\crefname{originallemma}{original lemma}{original lemmas}
\Crefname{originallemma}{Original lemma}{Original lemmas}
\makeatletter
\newlabel{topology-lemma-closed-morphism-quotient-topology}{{0AAU}{source}{}{sourcepointer-4}{}}
\newlabel{topology-lemma-closed-morphism-quotient-topology@cref}{{[originallemma][0][]0AAU}{[1][0][]source}{}{}{}}
\makeatother
\title{Universal property of the contraction of an exceptional curve of the first kind}
\author{The Stacks Project Authors}
\date{stacks-project:resolve.tex:local-snapshot-20261012; GNU FDL 1.2 or later}
\begin{document}
\maketitle
\noindent\textit{Editorial note.} All delivered mathematical text is the author's own, copied from the pinned Stacks Project chapter file resolve.tex; no statement or proof step was rewritten, reordered or removed.\par\medskip
\noindent\textit{Editorial note.} The counted proof is the author's complete argument: the identification of the contraction with the quotient space carrying the pushforward structure sheaf, the construction of the factorization on the level of ringed spaces through adjunction, and the verification of locality of the induced map on stalks at the image point.\par\medskip
\noindent\textit{Editorial note.} The label lemma-factor-through-contraction names the counted statement, and the proof environment immediately following it in the source is delivered verbatim as its proof.\par\medskip
\noindent\textit{Editorial note.} The counted statement is carried by the author's own lemma environment in the chapter source.\par\medskip
\noindent\textit{Editorial note.} Pre-open gate: python operations/agent-authoring/tools/precard\_gate.py --source resolve.tex --label lemma-factor-through-contraction --start 4538 --end 4589 --local-id 28090 returned RESULT=OK (L28090 free, occupancy FREE, public interval clear). Pre-assembly tag self-check: canonical tag resolved and not present in reference/public-items.jsonl.\par\medskip
\noindent\textit{Editorial note.} The statement and proof are copied from the pinned chapter file; the cited inputs are kept as external pointers with their canonical Stacks tags.\par\medskip
\noindent\textit{Allowed standard prerequisites (editorial).}\begin{itemize}
\item Properness of the contraction and submersiveness of a proper surjective map. Divisors, the lemma that a blowup is projective, Morphisms, the lemma that a locally projective morphism is proper, and Topology, the lemma that a closed surjective continuous map induces the quotient topology. Supplies the description of the underlying topological space of the contraction as the quotient identifying the points of the curve.
\item Computation of the pushforward of the structure sheaf along the contraction. Same chapter, the lemma computing the cohomology of the structure sheaf of the blowup, in particular the statement that the structure sheaf of the contraction maps isomorphically to the pushforward of the structure sheaf of the blowup. Supplies the identification of the structure sheaf of the contraction with the pushforward of the structure sheaf of the source.
\end{itemize}\medskip
\ifdefined\hypertarget\hypertarget{sourcepointer-1}{}\fi
\noindent\textit{Source reference.} Original lemma 02NS: sources/stacks-chow/tags.txt, source lines 3491--3491. Prerequisite pointer invoked at proof line(s) 4552; it names a standard fact of a sibling chapter that the counted argument uses, and every step of the counted proof built on it is retained verbatim in the delivered body.\par\medskip
\ifdefined\hypertarget\hypertarget{sourcepointer-2}{}\fi
\noindent\textit{Source reference.} Original reference 0AGS: sources/stacks-resolve/pinned-original-source.tex, source lines 786--786. Prerequisite pointer invoked at proof line(s) 4561; it names a standard fact of the same chapter that the counted argument uses, and every step of the counted proof built on it is retained verbatim in the delivered body.\par\medskip
\ifdefined\hypertarget\hypertarget{sourcepointer-3}{}\fi
\noindent\textit{Source reference.} Original lemma 01WC: sources/stacks-chow/tags.txt, source lines 2477--2477. Prerequisite pointer invoked at proof line(s) 4553; it names a standard fact of a sibling chapter that the counted argument uses, and every step of the counted proof built on it is retained verbatim in the delivered body.\par\medskip
\ifdefined\hypertarget\hypertarget{sourcepointer-4}{}\fi
\noindent\textit{Source reference.} Original lemma 0AAU: sources/stacks-chow/tags.txt, source lines 12997--12997. Prerequisite pointer invoked at proof line(s) 4556; it names a standard fact of a sibling chapter that the counted argument uses, and every step of the counted proof built on it is retained verbatim in the delivered body.\par\medskip
\begin{lemma}
\label{lemma-factor-through-contraction}
Let $X$ be a Noetherian scheme. Let $E \subset X$ be an
exceptional curve of the first kind. If a contraction $X \to X'$
of $E$ exists, then it has the following universal property:
for every morphism $\varphi : X \to Y$ such that $\varphi(E)$
is a point, there is a unique factorization
$X \to X' \to Y$ of $\varphi$.
\end{lemma}

\begin{proof}
Let $b : X \to X'$ be a contraction of $E$. As a topological space
$X'$ is the quotient of $X$ by the relation identifying all points
of $E$ to one point. Namely, $b$ is proper
(Divisors, Lemma \ref{divisors-lemma-blowing-up-projective} and
Morphisms, Lemma \ref{morphisms-lemma-locally-projective-proper})
and surjective, hence defines a submersive map of topological
spaces (Topology, Lemma
\ref{topology-lemma-closed-morphism-quotient-topology}).
On the other hand, the canonical map
$\mathcal{O}_{X'} \to b_*\mathcal{O}_X$ is an isomorphism. Namely,
this is clear over the complement of the image point $x \in X'$ of $E$
and on stalks at $x$ the map is an isomorphism by part (4) of
Lemma \ref{lemma-cohomology-of-blowup}.
Thus the pair $(X', \mathcal{O}_{X'})$ is constructed
from $X$ by taking the quotient as a topological space
and endowing this with $b_*\mathcal{O}_X$ as structure sheaf.

\medskip\noindent
Given $\varphi$ we can let $\varphi' : X' \to Y$ be the unique map of
topological spaces such that $\varphi = \varphi' \circ b$.
Then the map
$$
\varphi^\sharp : \varphi^{-1}\mathcal{O}_Y =
b^{-1}((\varphi')^{-1}\mathcal{O}_Y) \to \mathcal{O}_X
$$
is adjoint to a map
$$
(\varphi')^\sharp :
(\varphi')^{-1}\mathcal{O}_Y \to b_*\mathcal{O}_X = \mathcal{O}_{X'}
$$
Then $(\varphi', (\varphi')^\sharp)$ is a morphism of ringed spaces
from $X'$ to $Y$ such that we get the desired factorization. Since
$\varphi$ is a morphism of locally ringed spaces, it follows that
$\varphi'$ is too. Namely, the only thing to check is that the map
$\mathcal{O}_{Y, y} \to \mathcal{O}_{X', x}$ is local, where $y \in Y$
is the image of $E$ under $\varphi$. This is true because an element
$f \in \mathfrak m_y$ pulls back to a function on $X$ which is zero
in every point of $E$ hence the pull back of $f$ to $X'$ is a function
defined on a neighbourhood of $x$ in $X'$ with the same property.
Then it is clear that this function must vanish at $x$ as desired.
\end{proof}

\end{document}
